
\paragraph{\textbf{Proof of \SGARefLink{sga3:xiv:7:lemma3}{Lemma 3}.}}

The proof uses Coxeter transformations. Let
\(a_1,\ldots,a_n\) be a simple system of roots of \(R\), and let
\(r_i\) be the reflection corresponding to \(a_i\). Put
\[
  c=r_1\cdots r_n.
\]
Then \(c\in W\). Although \(c\) depends on the simple system and
on its ordering, its conjugacy class does not. It is called the
Coxeter element. We shall use, without proof, the fact that \(1\)
is not an eigenvalue of \(c\).

\paragraph{\textbf{Lemma 5.}}
\SGARefTarget{sga3:xiv:7:lemma5}
Put
\[
  \theta_i=r_n\cdots r_{i+1}(a_i).
\]
Then:
\begin{enumerate}
\item[\textup{\SGAInlineRefTarget{sga3:xiv:7:lemma5:item:a}(a)}] the \(\theta_i\) form a basis of the group
\(X\) generated by the roots;
\item[\textup{\SGAInlineRefTarget{sga3:xiv:7:lemma5:item:b}(b)}] \(\theta_i>0\) and \(c(\theta_i)<0\) for
every \(i\);
\item[\textup{\SGAInlineRefTarget{sga3:xiv:7:lemma5:item:c}(c)}] every root \(a\in R\) such that \(a>0\) and
\(c(a)<0\) is one of the \(\theta_i\);
\item[\textup{\SGAInlineRefTarget{sga3:xiv:7:lemma5:item:d}(d)}] \(R\) is the union of the orbits of the
\(\theta_i\) under the powers of \(c\).
\end{enumerate}

Indeed,
\[
  \theta_i=a_i+\sum_{j>i}m_{ij}a_j,\qquad m_{ij}\in\mathbb Z,
\]
which proves \textup{\SGARefLink{sga3:xiv:7:lemma5:item:a}{(a)}}. Assertions \textup{\SGARefLink{sga3:xiv:7:lemma5:item:b}{(b)}} and
\textup{\SGARefLink{sga3:xiv:7:lemma5:item:c}{(c)}} follow from the observation that \(r_i\) preserves
the sign of every root other than \(\pm a_i\), and reverses the
sign of \(\pm a_i\).

For \textup{\SGARefLink{sga3:xiv:7:lemma5:item:d}{(d)}}, an orbit of \(c\) cannot consist entirely of
positive roots, nor entirely of negative roots. Otherwise the
sum of the roots in that orbit would give a nonzero element of
\(X\) fixed by \(c\), contrary to the fact that \(1\) is not an
eigenvalue. Hence every orbit contains a root \(a>0\) for which
\(c(a)<0\), and \textup{\SGARefLink{sga3:xiv:7:lemma5:item:c}{(c)}} applies.

\paragraph{\textbf{Remark.}}
The proof was recalled only for the reader's convenience; one
could instead refer to the standard texts on Coxeter
transformations, for example Koszul, Séminaire Bourbaki
1959/1960, Exposé~191. They also show that all orbits of \(c\)
have the same number \(h\) of elements and each contains exactly
one \(\theta_i\). In particular,
\[
  h=\frac{\operatorname{Card}(R)}{n}.
\]

We return to \SGARefLink{sga3:xiv:7:lemma3}{Lemma 3} and distinguish three cases.

\medskip
\noindent\textup{\SGARefLink{sga3:xiv:7:lemma1}{(1)}}
Suppose that \(\varphi\in E/W\) is the identity. We must choose
\(f\in W\); take \(f=c\). \SGARefLink{sga3:xiv:7:lemma5}{Lemma 5} shows that this works.

\medskip
\noindent\textup{\SGARefLink{sga3:xiv:7:lemma2}{(2)}}
Suppose that \(R\) has type \(A_n\), where \(n\ge2\) is even, and
that \(\varphi\) is the unique nontrivial element of \(E/W\).
Then \(-1\notin W\), and we take
\[
  f=-c.
\]
The Coxeter element \(c\) has order \(h=n+1\), which is odd. If
\(a\in R\), \SGARefLink{sga3:xiv:7:lemma5}{Lemma 5} gives \(a=c^m\theta_i\) for suitable \(m,i\).
Adding \(h\) to \(m\) if necessary, we may suppose that \(m\) is
even. Then
\[
  a=(-c)^m\theta_i=f^m\theta_i.
\]
Thus the orbits of the \(\theta_i\) under \(f\) fill \(R\).
